% !TeX program = xelatex
% !TeX encoding = UTF-8
\documentclass{MathorCupmodeling}
\usepackage{zhlipsum,mwe}
\bianhao{MC3141592}
\tihao{（A/B/C/D）}
\timu{MathorCup 杯 \LaTeX{} 模板写作示例}
\keyword{content；content；content}
\begin{document}
	\begin{abstract}
		请使用 \TeX Live 2023，XeLaTeX 编译，请选用支持 UTF­-8 编码的编辑器。

		使用者需要有一定的 \LaTeX{} 的使用经验，至少要会使用常用宏包的一些功能，比如参考文献，数学公式，图片使用，列表环境等等。模板已经添加了常用的宏包，无需用户再额外添加。

		本模板
		\begin{itemize}
			\item 定义了几个宏 \lstinline|\def\ee{\mathrm{e}},\def\ii{\mathrm{i}},\def\leq{\leqslant},\def\geq{\geqslant}| 方便使用；
			\item 图片应放在 \lstinline|figure| 文件夹中；
			\item 定制了 matlab 和 python 代码环境，使用方法：\lstinline|\begin{matlab} content \end{matlab}| 和 \lstinline|\begin{python} content \end{python}|；
			\item 加载了 \lstinline|cleveref| 宏包，使用方法：\lstinline|\cref{label}|。
		\end{itemize}
		其它的就是跟普通的 \lstinline|ctexart| 使用方法一样。
	
		欢迎到 \url{https://ask.latexstudio.net/} 或者 QQ 群 200392395  提问，注意提供 MWE。
		
		
		
字数尽量控制在 300$\sim $600之间，需控制在一页；摘要中必须将具体方法、模型和所得结果写出来；摘要要求“总分总”，段开头可用“针对问题1，针对问题2，针对问题3..”或者“首先，然后，其次，最后”等词语进行有逻辑的论述。摘要是重中之重，必须认真严谨！
	\end{abstract}
	\tableofcontents\newpage

\section{问题重述}
\subsection{引言}
可以自己组织词句对问题进行描述，主要数据可以直接复制，对所提出的问题部分基本原样复制。篇幅建议不要超过一页。大部分文字提炼自原题。

\subsection{要解决的具体问题}
\begin{enumerate}
  \item \textcolor{red}{问题一：问题一的重述，重述语言简洁明了，突出模型中第一个要解决的问题，突出核心。}
  \item \textcolor{red}{问题二：问题二的重述。}
  \item \textcolor{red}{问题三：问题三的重述。}
  \item \textcolor{red}{问题四：问题四的重述。}
\end{enumerate}

\section{问题分析}

\textcolor{red}{主要是表达对题目的理解，特别是对附件的数据进行必要分析、描述（一般都有数据附件），这是需要提到分析数据的方法、理由。如果有多个小问题，可以对每个小问题进行分别分析。问题分析中不给出结果，结果在摘要中给出。
（假设有2个问题）}

\subsection{问题一的分析}

 \textcolor{red}{对问题1研究的意义的分析。
问题1属于$\cdots\cdots$数学问题，对于解决此类问题一般数学方法的分析。
对附件中所给数据特点的分析。
对问题1所要求的结果进行分析。
由于以上原因，我们可以将首先建立一个$\cdots\cdots$的数学模型I,然后将建立一个$\cdots\cdots$的模型II,$\cdots\cdots$对结果分别进行预测，并将结果进行比较.
}


\subsection{问题二的分析}
 \textcolor{red}{对问题2研究的意义的分析。
问题2属于$\cdots\cdots$数学问题，对于解决此类问题一般数学方法的分析。
对附件中所给数据特点的分析。
对问题2所要求的结果进行分析。
由于以上原因，我们可以将首先建立一个$\cdots\cdots$的数学模型I,然后将建立一个$\cdots\cdots$的模型II,$\cdots\cdots$对结果分别进行预测，并将结果进行比较.
}

 $\cdots\cdots$


\section{模型假设}
\begin{enumerate}
  \item 模型的假设要结合整个模型的建立作出的一个合理的假设，不能过于理想化，要尽量切合实际问题的处理来做出相应的合理的假设；
  \item 模型假设二；
  \item 模型假设三；
\end{enumerate}

\section{名词解释与符号说明}

\textcolor{red}{一般都会有符号解释和说明，对于一些装有的专有名词解释，需要的时候就需要对其进行解释与说明，我们以下面几个例子为例。}

\subsection{名词解释与说明}
\begin{enumerate}
\item \textbf{理论通行能力：}理论通行能力是指每一条车道~(或每一条道路) 在单位时间内
能够通过的最大交通量。


\item \textbf{修正通行能力：}在具体条件下，通过修正系数对理论通行能力修正后得到的单
位时间内所能通过的最大交通量。

\end{enumerate}
\subsection{主要符号与说明}

%tab1
\begin{table}[h!]
  \centering
  \small
  \begin{tabular}{p{60pt}<{\centering}|p{60pt}<{\centering}p{180pt}<{\raggedright}}
   \hline
    序号 & 符号 & 符号说明 \\
   \hline
    1 & $\nu$ & 行车速度（km/h） \\
    2 & t$_{\min}$ & 车头最小时距（s） \\
    3 & $J_{\rm a}$ & 车头最小间隔（m） \\
    4 & $J_{\rm z}$ & 车辆平均长度（m） \\
    5 & $J_{\gamma}$ & 车辆的制动距离（m） \\
    6 & $J_{\max}$ & 司机在反应时间内车辆行驶的距离（m） \\
    7 & $A_{\max}$ & 最大交通量 \\
    8 & $\alpha_{1}$ & 车道数修正系数 \\
    9 & $\alpha_{2}$ & 车道宽度和侧向净宽修正系数 \\
    10 & $\alpha_{3}$ & 大型车修正系数 \\
    11 & $\alpha_{4}$ & 驾驶员技术水平修正系数 \\
    12 & $K_{j}$ & 阻塞密度 \\
    13 & $\nu_{f}$ & 自由车速 \\
    $\cdots$ & $\cdots$\\
    \hline
  \end{tabular}
  %\caption{符号与说明}
  \label{symbol}
\end{table}

\section{模型的建立与求解}

数据的预处理：
1. $\cdots\cdots$数据全部缺失，不予考虑。
2. 对数据测试的特点，如周期等进行分析。
3. $\cdots\cdots$数据残缺，根据数据挖掘等理论根据$\cdots\cdots$变化趋势进行补充。
4. 对数据特点（后面将会用到的特征）进行提取。
   用$\cdots\cdots$软件聚类分析和各个不同问题的需要，采得$\cdots\cdots$组采样，每组5-8个采样值。将采样所对应的特征值进行列表或图示。
根据数据特点，对总体和个体的特点进行比较，以表格或图示方式显示。


\subsection{问题一的分析和求解}

\subsubsection{***模型的建立}

模型建立的内容要点如下：

模型的主要类别：

几种常见的建模目的：

建模过程常见的几个要点：

模型的基本要求：

模型选择要点：

加分项（能在规定时间内做完后还有足够时间的再考虑加分项）：

1、鼓励创新。在能解决问题的基础上，对经典模型进行改进，欣赏独树一帜、有创新性的模型，但要合理。

2、对于同一问题使用两个或以上合理模型进行求解。避免出现单纯罗列模型，又不做对比和评价的现象。
 
 

	\section{模型的建立与求解}
	\zhlipsum*[10]
	\subsection{模型的准备}
	\zhlipsum*[11]
	\subsection{模型一的建立}
	效果见\cref{fig:1}。
	\begin{figure}[htbp]
		\centering
		\includegraphics{example-image-plain.pdf}
		\caption{content}\label{fig:1}
	\end{figure}
	\subsection{模型二的建立}
	结果见\cref{tab:1}。
	\begin{table}[htbp]
		\centering
		\caption{content}\label{tab:1}
		\begin{tabularx}{0.7\textwidth}{c@{\hspace{1pc}}|@{\hspace{2pc}}X}
		\Xhline{0.08em}
		符号 & \multicolumn{1}{c}{符号说明}\\
		\Xhline{0.05em}
		$\delta$ & 赤纬角\\
		$\beta$ & 经度\\
		$\alpha$ & 纬度\\
		$r$ & 地球半径\\
		$\gamma$ & 太阳光与杆所成的夹角\\
		$l$ & 杆的长度\\
		$l_{y}$ & 杆的影子长度\\
		$\vec{x}_{1},\vec{y}_{1},\vec{z}_{1}$ & 由杆的位置所生成的切平面的正交基\\
		$\vec{\hat{x}}_{1},\vec{\hat{y}}_{1},\vec{\hat{z}}_{1}$ & 由杆的位置所生成的切平面的单位正交基\\
		$\theta$ & 影子与北方的夹角\\
		$l_{y}(i)$ & 编号为 $i$ 的数据对应的影子长度\\
		$\theta_{i}$ & 编号为 $i$ 的数据对应的影子角度\\			\Xhline{0.08em}
		\end{tabularx}
	\end{table}
	\subsection{模型三的建立}
	\zhlipsum*[14]

	\section{模型的检验}
	\zhlipsum*[15]

	\section{模型的评价与改进}
	\zhlipsum*[16]
	\subsection{模型的优点}
	\zhlipsum*[17]
	\subsection{模型的缺点}
	\zhlipsum*[18]
	\subsection{模型的改进}
	\zhlipsum*[19]\cite{label}

	\phantomsection
	\addcontentsline{toc}{section}{参考文献}
	\begin{thebibliography}{99}
	\bibitem{label}https://www.latexstudio.net/
	\end{thebibliography}

	\newpage
	\appendix
	\ctexset{section={
		format={\zihao{-4}\heiti\raggedright}
	}}
	\begin{center}
		\heiti\zihao{4} 附\hspace{1pc}录
	\end{center}
	\section{问题一的 MATLAB 代码}
	\begin{matlab}
clc,clear
%第七题
R71 = 1;
R72 = 2;
T7 = 1;
K7 = 1;
N7=10^5;
G71=tf(R72,R71);
G72=tf(1,[T7 1]);
G73=tf(1,[T7 0]);
G74=tf([N7*T7 0],[T7 N7]);
G75=tf([N7*K7*T7 N7*K7],[T7 N7]);
G76=tf([K7*T7 K7],[T7 0]);

subplot(2,3,1)
step(G71)
xlabel('$t$','interpreter','latex', 'FontSize', 12);
ylabel('$y$','interpreter','latex', 'FontSize', 12);
title('比例环节');

subplot(2,3,2)
step(G72)
xlabel('$t$','interpreter','latex', 'FontSize', 12);
ylabel('$y$','interpreter','latex', 'FontSize', 12);
title('惯性环节');

subplot(2,3,3)
step(G73)
xlabel('$t$','interpreter','latex', 'FontSize', 12);
ylabel('$y$','interpreter','latex', 'FontSize', 12);
title('积分环节');

subplot(2,3,4)
step(G74,10^-3)
xlabel('$t$','interpreter','latex', 'FontSize', 12);
ylabel('$y$','interpreter','latex', 'FontSize', 12);
title('微分环节');

subplot(2,3,5)
step(G75,10^-3)
xlabel('$t$','interpreter','latex', 'FontSize', 12);
ylabel('$y$','interpreter','latex', 'FontSize', 12);
title('比例微分环节');

subplot(2,3,6)
step(G76)
xlabel('$t$','interpreter','latex', 'FontSize', 12);
ylabel('$y$','interpreter','latex', 'FontSize', 12);
title('比例积分环节');
	\end{matlab}
\end{document} 